% !TEX root = ../CerdaMS.tex
\newcounter{theoremtmp}

\section{Proofs of Section \ref{sec:MemTrees}}

\setcounter{theoremtmp}{\arabic{theorem}}
\setcounter{theorem}{\arabic{lem:bijAppsetfinmemtrees}}
\begin{lemma}[recalled from \cpageref{lem:bijAppsetfinmemtrees}]
	There is a bijection between $\Appset$ and the set of finite Ohana 
	trees of \lamI-terms.
\end{lemma}
\setcounter{theorem}{\arabic{theoremtmp}}

	\begin{proof}
	Given $A\in\Appset$, define a \lamI-term $M_A$ 
	by structural induction on $A$:
	\begin{itemize}
	\item if $A = \bot_{\set{x_1,\dots,x_n}}$ then 
		$M_A\eqdef \Om x_1\cdots x_n$.
	\item if $A = \lam\vec x.yA_1\cdots A_k$ then 
		$M_A \eqdef \lam\vec x.y(M_{A_1})\cdots (M_{A_k})$.
	\end{itemize}
	Define a map $\iota$ associating to $A$ the Ohana tree $\MT{M_A}$, 
	i.e.\ $\iota(A) \eqdef \MT{M_A}$.
	
	Conversely, let $M\in\LamI$ be such that $T\eqdef\MT{M}$ is finite. 
	Since $T$ is finite, we can construct an approximant $\iota^-(T)$ 
	by structural induction on $T$. There are two cases:
	\begin{itemize}
	\item $T = \bot_X$, with $X = \fv(M)$. 
		In this case, define $\iota^-(T) \eqdef \bot_X$.
	\item If $T = \lam\vec x.y\cdot^{X_1}T_1\cdots^{X_k} T_k$ then 
		$M\redd[h] \lam\vec x.yM_1\cdots M_k$ with $T_i = \MT{M_i}$ 
		and $X = \fv(M_i)$, for all $i\,(1\le i\le k)$. 
		In this case, define $\iota^-(T) = 
		\lam\vec x.y(\iota^-(T_1))\cdots (\iota^-(T_k))$.
	\end{itemize}
	A straightforward induction shows $\iota(\iota^-(T)) = T$ and 
	$\iota^-(\iota(A)) = A$.
	\end{proof}

\setcounter{theoremtmp}{\arabic{theorem}}
\setcounter{theorem}{\arabic{lem:DirApp}}
\begin{lemma}[recalled from \cpageref{lem:DirApp}] \mbox{}
\bsub
\item
	For every $M,N\in\LamI$, $M\bred N$ entails  $\dirapp{M} \MTle 
	\dirapp{N}$.
\item
	For all $M\in\LamI(X)$, $\Appof{M}$ is an ideal of the poset 
	$(\Appset(X),\sqsubseteq)$. More precisely:
	\begin{enumerate}
	\item
		non-emptiness: $\bot_X\in \Appof{M}$;
	\item
		downward closure: $A\in \Appof{M}$ and $A'\MTle A$ imply 
		$A'\in\Appof{M}$;
	\item
		directedness: $\forall A_1,A_2\in\Appof{M},\exists A_3\in\Appof{M}$ 
		such that $A_1\MTle A_3 \sqsupseteq A_2$.
	\end{enumerate}
\esub
\end{lemma}
\setcounter{theorem}{\arabic{theoremtmp}}

	\begin{proof}
	\bsub
	\item By induction on a derivation of $M\bred N$. By \cref{rem:hnfs}, 
		either $M$ has a head redex or it has a head variable.
		If $M = \lam \vec x.(\lam y.M')M_0M_1\cdots M_k$ then 
		$\dirapp{M} = \bot_{\fv(M)}$. 
		By \cref{rem:aboutLambdaI}\eqref{rem:aboutLambdaI2}, 
		$M\bred N$ entails $\fv(N) = \fv(M)$, 
		whence $\dirapp{N}\in\Appset(\fv(M))$.
		This case follows since $\bot_{\fv(M)}$ is the bottom of 
		$\Appset(\fv(M))$.
		
		If $M = \lam \vec x.yM_1\cdots M_k$ then $\dirapp{M} = \lam \vec 
		x.y\,\dirapp{M_1}\cdots \dirapp{M_k}$.
		The $\beta$-reduction must occur in some $M_i\bred N_i$, whence 
		$\dirapp{N} = \lam \vec x.y\,\dirapp{M_1}\cdots \dirapp{N_i}\cdots 
		\dirapp{M_k}$. We conclude since $\dirapp{M_i}\MTle\dirapp{N_i}$ 
		holds, from the induction hypothesis.
			
		\item To prove that $\Appof{M}$ is an ideal, we need to check three 
		conditions.
		\begin{enumerate}
		\item
			non-emptiness: since $\bot_{X}\sqsubseteq \dirapp{M}$ for all 
			$M\in\Appset(X)$, we get $\bot_{X} \in \Appof{M}$.
		\item
			downward closure: by definition of $\Appof{M}$.
		\item
			directedness: let us take $A_1,A_2\in\Appof{M}$ and prove that 
			$A_1\sqcup A_2\in\Appof{M}$.
			We proceed by structural induction on, say, $A_1$.
			
			Case $A_1 = \bot_X$. Then $A_2\in\Appset(X)$ and $A_1\sqcup A_2 
			= A_2\in\Appof{M}$.
			
			Case $A_1 = \lam x_1\ldots x_n.y\,A'_1\cdots A'_k$. If $A_2 = 
			\bot_X$ then proceed as before. Otherwise $A_1\in\Appof{M}$ 
			because there is a reduction $M\redd[\beta] \lam x_1\ldots 
			x_n.y\,M'_1\cdots M'_k=:M_1$ with $A'_i\in\Appof{M'_i}$.
			Similarly, $A_2\in\Appof{M}$ because there is a reduction 
			$M\redd[\beta] \lam x_1\ldots x_{n'}.y\,M''_1\cdots 
			M''_{k'}=:M_2$ with $A_2 = \lam x_1\ldots x_{n'}.y\,A''_1\cdots 
			A''_{k'} $ and $A''_j\in\Appof{M''_j}$.
			By confluence of $\bred$, $n = n'$ and $k = k'$, and every 
			$M'_i,M''_i$ have a common reduct $M'''_i$.
			By \cref{lem:DirApp}\eqref{lem:DirApp1}, 
			$A'_i,A''_i\in\Appof{M'''_i}$, so by HI $A'_i\sqcup 
			A''_i\in\Appof{M'''_i}$.
			Conclude by taking $A_3 \eqdef \lam x_1\ldots x_n.y\,(A'_1\sqcup 
			A''_1)\cdots (A'_k\sqcup A''_k)$.\qedhere
		\end{enumerate}
	\esub
	\end{proof}

\setcounter{theoremtmp}{\arabic{theorem}}
\setcounter{theorem}{\arabic{prop:InvariantUnderBeta}}
\begin{proposition}[recalled from \cpageref{prop:InvariantUnderBeta}]
	Let $M,N\in\LamI(X)$. If $M\red[\beta] N$, then $\Appof{M} = \Appof{N}$.
\end{proposition}
\setcounter{theorem}{\arabic{theoremtmp}}

	\begin{proof}
	To prove $\Appof{M} = \Appof{N}$, we show the two inclusions.
	
	($\subseteq$)
	Take $A\in\Appof{M}$, then there exists a reduction $M\redd[\beta] M'$, 
	such that $A\MTle\dirapp{M'}$. We proceed by induction on $A$.
	
	Case $A = \bot_X$. By Remark~\ref{rem:appof-M-fv-M}, since $\bot_X\MTle 
	\dirapp{N}$, for all $N\in\LamI(X)$.
	
	Case $A = \lam\vec x.yA_1\cdots A_k$. 
	Then $M' = \lam\vec x.yM'_1\cdots M'_k$ with $A_i\MTle\dirapp{M'_i}$, 
	for all $1\le i \le k$.
	Since $M\red[\beta] N$, by confluence, $N$ and $M'$ have a common reduct 
	$N' = \lam\vec x.yN'_1\cdots N'_k$ with $M'_i\redd[\beta] N'_i$. By 
	\cref{lem:DirApp}\eqref{lem:DirApp1} we have 
	$\dirapp{M'_i}\MTle\dirapp{N'_i}$ and, by transitivity, 
	$A_i\MTle\dirapp{N'_i}$.
	We found a reduction $N\redd[\beta] N'$ such that $A\MTle \dirapp{N'}$, 
	whence $A\in\Appof{N}$.
	
	($\supseteq$) Straightforward, since every reduct of $N$ is also a 
	reduct of $M$.
	\end{proof}


For the proof of \cref{thm:MTbijApp}, we introduce some definitions.
%\comgm{Alexis: Enjoy}
%\begin{definition}
%	Given an Ohana tree $T = \MT{M}$, one defines
%	$\fv(T)$ to be $\fv(M)$.
%\end{definition}
%The following proposition is immediate:\medskip
First, given an Ohana tree $T = \MT{M}$, one defines
$\fv(T)$ to be $\fv(M)$.
%\begin{proposition}
If 
$T=\lam x_1\ldots x_n.y\cdot^{X_1}T_1\cdots^{X_k}T_k$.
%$begin{tikzpicture}
%	\node (MTM) {$T = \lam x_1\dots x_n.y~$};
%	\node at (1.65,-1) {${T_1}\quad\cdots\quad {T_k}$};
%	\draw (0.6,-20pt) to[pos=0.4]node{$\qquad\ X_1$} (1.45,-5pt);
%	\draw (2.8,-20pt) to[pos=0.35]node{$\quad\ldots\quad X_k$} (1.7,-5pt);
%\end{tikzpicture}
then $\fv(T) = \bigcup_{i=1}^k X_i \cup \{y\} \setminus \{x_1,\dots,x_n\}\}\}$.
%\end{proposition}
Moreover, in the following, we will consider the class of head contexts, defined inductively by the following grammar:
${\sf H} ::= [] \mid \lam x_1\ldots x_n.y{\sf H}_1\cdots {\sf H}_k.$



Then, in order to relate Ohana trees with approximants, 
we extend the order on approximants to Ohana trees and define 
the set of finite subtrees of an Ohana tree:

\begin{definition} Let ${\cal M}=\set{\MT{M}}[M\in\LamI]$.
	We define $\mathord{\MTle}\subseteq {\cal M}^2$ as 
	the least partial order on Ohana trees closed under the following rules:
	\[
	\infer{\bot_X\MTle T}{\fv(T)=X}\qquad
	\infer{\lam x_1\ldots x_n.y\cdot^{X_1}T_1\cdots^{X_k}T_k
		\MTle \lam x_1\ldots x_n.y\cdot^{X_1}T'_1\cdots^{X_k}T'_k}
		{T_i\MTle T'_i&\textrm{for all }1\le i \le k}
	\]
\end{definition}


\begin{definition}
	Given an Ohana tree $T = \MT{M}$, one defines
	$T^\star$ to be $\set{T'}[T'\MTle T \text{ and $T'$ is finite}]$.
\end{definition}

\newcommand\pos[1]{{\sf Pos}(#1)}

\begin{definition}
	We define inductively the set of positions of an approximant $A$, 
	$\pos{A}$ as a subset of $\Nat^\star$ such that:
	\begin{itemize}
		\item If $A=\bot_X$, $\pos{A}=\epsilon$;
		\item If $A= \lam x_1\ldots x_n.yA_1\cdots A_k$, 
		$\pos{A}=\set{i\cdot p}[1\leq i \leq k \mathop{\&} p\in\pos{A_i}]$.
	\end{itemize}
	Sets of positions are lifted to sets of approximants as expected 
	and one can similarly define $\pos{M} = \pos{\Appof{M}}$. 
\end{definition}

\begin{definition}
	Given an approximant $A$ and $p\in \pos{A}$, one defines 
	$A(p)$ by induction on the length of $p$ as follows:
	\begin{itemize}
		\item If $p=\epsilon$, then if $A=\bot_X$, ${A}(p)=\bot_X$,
		otherwise, $A= \lam x_1\ldots x_n.yA_1\cdots A_k$ and 
		$A(p)=\lam x_1\ldots x_n.y$;
		\item If $p=i\cdot p'$, then $A= \lam x_1\ldots x_n.yA_1\cdots A_k$, 
		with $1\leq i \leq k$ and $A(p)=A_i(p')$.
	\end{itemize}
\end{definition}

%\begin{theorem}
%		For all $M\in\LamI$, we have an isomorphism $\MT{M} \cong \Appof{M}$.
%\end{theorem}

\setcounter{theoremtmp}{\arabic{theorem}}
\setcounter{theorem}{\arabic{thm:MTbijApp}}
\begin{theorem}[recalled from \cpageref{thm:MTbijApp}]
	There is a bijection $\mtapp{\cdot}$ between the set of Ohana trees 
	and the set of approximants of \lamI-terms satisfying 
	$\mtapp{\MT{M}} = \Appof{M}$, for all $M\in\LamI$.
%	For all $M\in\LamI$, there is an isomorphism $\MT{M} \cong \Appof{M}$.
\end{theorem}
\setcounter{theorem}{\arabic{theoremtmp}}

%\comasi{Suggestion of reformulation of the theorem statement:}

\begin{proof}[Proof of \cref{thm:MTbijApp}] 
	Let us write $\mathbb{T} = \set{\MT{M}}[M\in\LamI]$ and 
	$\mathbb{A} = \set{\Appof{M}}[M\in\LamI]$ and 
	let us build a bijection ${\cal A}: \mathbb{T} \rightarrow \mathbb{A}$ 
	such that for any $M\in\LamI$, $\mtapp{\MT{M}} = \Appof{M}$.
	Let us define ${\cal A}$.
	%	From $\MT{M}$ to $\Appof{M}$.
	Let $M\in\LamI$ and $T$ its Ohana tree. 
	Since $T^\star$ contains only finite "subtrees" of $T$, 
	which are all Ohana trees of $\lamI$-terms, 
	one can consider a set of approximants $\mtapp T$ obtained 
	as the direct image of $T^\star$ by the mapping $\iota^-$
	defined in the proof of Lemma~\ref{lem:bijAppsetfinmemtrees}. 
	%Since $\iota^-$ is a bijection and $(\cdot)^\star$ is injective, 
	%their composition is so and i
	
	In other words, we define $\mtapp{T} \eqdef \iota^-(T^\star)$. 
	We now prove that $\mtapp T = \Appof{M}$.
	
	Let us prove that $\mtapp T \subseteq \Appof{M}$.
	Let $T'\in T^\star$, then there exist 
	$ M_1\in\LamI(X_1), \dots, M_n\in \LamI(X_n)$ and 
	a head context ${\sf H}$ such that 
	${\sf H}[M_1,\dots, M_n] \in\LamI$ 
	with $M\redd[\beta] {\sf H}[M_1,\dots, M_n]$ 
	and $T = \MT{{\sf H}[\Om\vec{X_1},\dots, \Om\vec{X_n}]}$.
	Therefore $\iota^-(T') = {\sf H}[\bot_{X_1}, \dots, \bot_{X_n}]$ 
	and since $\bot_{X_i}\MTle \dirapp{M_i}$, $\iota^-(T') \in\Appof{M}$ 
	which proves the inclusion. 
	
	For the converse inclusion, 
	%let us consider $A\in \Appof{M}$. One can associate to $A$ 
	%the length $l_A$ of the shortest standard reduction from $M$ 
	%that reaches a term $N$ such that $A \MTle \dirapp{N}$. 
	we prove by induction on the structure of $A$ that, for any $M\in\LamI$ 
	and $A\in\Appof{M}$, $A\in\mtapp{\MT M}$:
	\begin{itemize}
		\item If $A= \bot_X$, then $X = \fv(M)$ and 
		$A\in \iota^-(\MT{M}^\star)$;
		\item Otherwise, $A = \lam x_1\ldots x_n.yA_1\cdots A_k$. 
		In such a case, we have $M\red[h] \lam x_1\ldots x_n.yM_1\cdots M_k$ 
		with $M_i\in\LamI(X_i)$ and $A_i\in\Appof{M_i}$ 
		for $1\leq i \leq k$. 
		By induction hypothesis we have some $T_i \in\MT{M_i}^\star$ 
		such that $A_i\in \iota^-(T_i)$ for all $i\,(1\leq i \leq k)$. 
		Since $\lam x_1\ldots x_n.y \cdot^{X_1} T_1 \cdots^{X_k} T_k 
		\in\MT{M}^\star$, we indeed have that $A \in \mtapp{\MT{M}}$.
	\end{itemize}
	
	%	From $\Appof{M}$ to $\MT{M}$. 
	Let us define ${\cal T}$, the inverse of ${\cal A}$.
	Let $A, B\in \Appof{M}$. 
	%There exist $M_A, M_B\in \Lambda_I$ such that $M\redd[\beta] M_A, M_B$, 
	%$A\MTle \dirapp{M_A}$ and  $B\MTle \dirapp{M_B}$. 
	%Let $N$ be such that $M_A, M_B \redd[\beta] N$, we have by 
	By directedness of $\Appof{M}$ 
	(Lemma~\ref{lem:DirApp}(\ref{lem:DirApp2})), 
	there exists $C\in\Appof{M}$ such  that $A,B\MTle C$.
	From this, it follows that at any position where 
	two approximants of $M$ are defined and different from $\bot_X$, 
	they agree. 
	We  define a (potentially) infinite tree $\appmt{\Appof{M}}$ 
	from $\Appof{M}$ as a pair of partial functions $(v,e)$, 
	respectively labelling the edges and vertices of the tree, 
	defined on $\pos{\Appof{M}}$ as follows: 
	\begin{itemize}
		\item for every position $p \in \pos{\Appof{M}}$, 
		either for every $A\in \Appof{M}$ such that $p\in\pos{A}$, 
		$A(p)=\bot_X$, in which case $v(p) = \bot_X$, 
		or there is some $A\in \Appof{M}$ such that $p\in\pos{A}$, 
		$A(p)=\lam x_1\ldots x_n.y$ in which case 
		$v(p)\eqdef\lam x_1\ldots x_n.y$.
		\item for every non-empty position $p \in \pos{\Appof{M}}$,
		let $A\in \Appof{M}$ be such that $A(p)=\bot_X$ and 
		set $e(p)\eqdef X$.
	\end{itemize}
	A simple coinduction gives us 
	%$\appmt{\Appof{M}}=\MT{M}$.
	$\appmt{\mtapp{\MT{M}}}=\MT{M}$ for any $M\in\LamI$.
\end{proof}

\section{Proofs of Section \ref{sec:resource}}


\subsection{Substitution lemma (Lemma \ref{lem:resource-substitution})}

% \comrci{Alternative version: get rid of the \ensuremath{\lambda}I-hypotheses
%	in the substitution lemma.
%	We thought it was bad, but couldn't find again the reason why...}

% In the non-\ensuremath{\lambda}I version, no need for this lemma.
\begin{lemma} \label{lem:rsubst-x-notin-fv}
	Given $\rt s \in \rterms**(X)$, $\rb u \in \rterms*(Y)$
	and $x \notin X$,
	\[	\rt s \rsubst u = \left\{ \begin{array}{ll}
		\rt s, & \text{if $\rb u = 1_Y$,} \\
		\esum X, & \text{otherwise.}
		\end{array} \right. \]
\end{lemma}

	\begin{proof}
	Straightforward induction on $\rt s$.
	\end{proof}

\setcounter{theoremtmp}{\arabic{theorem}}
\setcounter{theorem}{\arabic{lem:resource-substitution}}
\begin{lemma}[recalled from \cpageref{lem:resource-substitution}]
	Given $\rt s \in \rterms**(X)$,
	$\rb u \in \rterms*(Y)$,
	$\rb v \in \rterms*(Z)$,
	$x \in X \setminus Z$ and $y \in X \cup Y$,
%	and $x,y \in \Var$ such that $x \notin Z$,
	\[	\rt s \rsubst u \rsubst v [y] =
		\sum_{\rb v = \rb v' \cdot \rb v''} \rt s \rsubst* {\rb v'} [y]
		\rsubst*{ \rb u \rsubst* {\rb v''}[y] }. \]
\end{lemma}
\setcounter{theorem}{\arabic{theoremtmp}}
	
	\begin{proof}
	By induction on $\rt s$.
	\begin{itemize}
	\item Case $\rt s = x$.
		(It is the only possible case of a variable
		because of the hypothesis $x \in X$.)
		We have
		\[	\sum_{\rb v = \rb v' \cdot \rb v''} x \rsubst* {\rb v'} [y]
				\rsubst*{ \rb u \rsubst* {\rb v''}[y] }
			= x \rsubst* {1_Z}[y] \rsubst*{\rb u \rsubst v[y]}
			= x \rsubst*{\rb u \rsubst v[y]}, \]
		hence if $\rb u = [\rt u]$, then
		both sides of the equality are equal to $\rt u \rsubst v[y]$,
		otherwise they are equal to $\esum{Y \msubst Z [y]}$.
%	\item Case $\rt s = y$. We have
%		\[	\sum_{\rb v = \rb v' \cdot \rb v''} y \rsubst* {\rb v'} [y]
%				\rsubst*{ \rb u \rsubst* {\rb v''}[y] }
%			= \sum_{\rb v = [\rt v] \cdot \rb v''} y \rsubst* {[\rt v]} [y]
%				\rsubst*{ \rb u \rsubst* {\rb v''}[y] }
%			= \sum_{\rb v = [\rt v] \cdot \rb v''}
%				\rt v \rsubst*{ \rb u \rsubst* {\rb v''}[y] }. \]
%		However, $x \notin Z$, therefore
%		if $\rb u = 1_X$ and $\rb v = [v]$,
%		both sides of the equality are equal to $v$,
%		otherwise they are equal to $\esum Y$.
%	\item Case $\rt s = z \notin \{x,y\}$.
%		If $\rb u = 1_X$ and $\rb v = 1_Y$,
%		both sides of the equality are equal to $z$.
%		Otherwise, they are equal to $\esum{\set z}$.
	\item Case $\rt s = \ensuremath{\lambda}y.\rt s'$.
		Immediate by the induction hypothesis on $s'$.
	\item Case $\rt s = \rt s' \rb s''$.
		We have $X = \fv(\rt s) = \fv(\rt s') \cup \fv(\rb s'')$.
		There are several possible cases.
		\begin{itemize}
		\item If $x \in \fv(\rt s')$ and $x \in \fv(\rb s'')$,
			then it is possible to apply the induction hypothesis
			in both subterms.
			By the definition of resource substitution,
			\begin{align*}
			\rt s \rsubst u \rsubst v [y]
			& = \adjustlimits \sum_{\rb u = \rb u_l \cdot \rb u_r}
				\sum_{\rb v = \rb v_l \cdot \rb v_r}
				\rt s' \rsubst* {\rb u_l} \rsubst* {\rb v_l} [y]
				\rb s'' \rsubst* {\rb u_r} \rsubst* {\rb v_r} [y] \\
			& = \adjustlimits \sum_{\rb u = \rb u_l \cdot \rb u_r}
				\sum_{\rb v = (\rb v'_l \cdot \rb v''_l) 
					\cdot (\rb v'_r \cdot \rb v''_r)}
				\rt s' \rsubst* {\rb v'_l} [y]
					\rsubst*{ \rb u_l \rsubst* {\rb v''_l}[y] }
				\rb s'' \rsubst* {\rb v'_r} [y]
					\rsubst*{ \rb u_r \rsubst* {\rb v''_r}[y] } \\
			\intertext{by the induction hypotheses 
			on $\rt s'$ and $\rb s''$,}
			& = \sum_{\rb v = \rb v' \cdot \rb v''} \rt s 
				\rsubst* {\rb v'} [y]
				\rsubst*{ \rb u \rsubst* {\rb v''}[y] }
			\end{align*}
			by permuting the sums and re-arranging the indices
			(using associativity and commutativity of the multiset union).
		\item If $x \in \fv(\rt s)$ and $x \notin \fv(\rb s'')$, then again
			\begin{align*}
			\rt s \rsubst u \rsubst v [y]
			& = \adjustlimits \sum_{\rb u = \rb u_l \cdot \rb u_r}
				\sum_{\rb v = \rb v_l \cdot \rb v_r}
				\rt s' \rsubst* {\rb u_l} \rsubst* {\rb v_l} [y]
				\rb s'' \rsubst* {\rb u_r} \rsubst* {\rb v_r} [y] \\
			& = \sum_{\rb v = \rb v_l \cdot \rb v_r}
				\rt s' \rsubst* {\rb u} \rsubst* {\rb v_l} [y]
				\rb s'' \rsubst* {1_Y} \rsubst* {\rb v_r} [y] \\
				& \qquad + \sum_{\substack{
					\rb u = \rb u_l \cdot \rb u_r \\ \rb u_r \neq 1_Y }}
				\sum_{\rb v = \rb v_l \cdot \rb v_r}
				\rt s' \rsubst* {\rb u_l} \rsubst* {\rb v_l} [y]
				\rb s'' \rsubst* {\rb u_r} \rsubst* {\rb v_r} [y] \\
			& = \sum_{\rb v = \rb v_l \cdot \rb v_r}
				\rt s' \rsubst* {\rb u} \rsubst* {\rb v_l} [y]
				\rb s'' \rsubst* {\rb v_r} [y]
				+ \esum{X \msubst Y \msubst Z[y]} \\
			\shortintertext{by \cref{lem:rsubst-x-notin-fv},}
			& = \sum_{\rb v = (\rb v'_l \cdot \rb v''_l) \cdot \rb v_r} 
				\rt s' \rsubst* {\rb v'_l} [y]
				\rsubst*{ \rb u \rsubst* {\rb v''_l}[y] }
				\rb s'' \rsubst* {\rb v_r} [y] \\
			\shortintertext{by the induction hypothesis on $\rt s'$,}
			& = \sum_{\rb v = (\rb v'_l \cdot \rb v_r) \cdot \rb v''_l}
				\rt s \rsubst* {\rb v'_l \cdot \rb v_r} [y]
				\rsubst* { \rb u \rsubst* {\rb v''_l}[y] }
			\end{align*}
			by re-arranging the indices,
			and using again \cref{lem:rsubst-x-notin-fv}
			together with the hypotheses
			$x \notin Z$ and $x \notin \fv(\rb s'')$.
			This is exactly the desired result.
		\item The case $x \notin \fv(\rt s)$ and $x \in \fv(\rb s'')$
			is symmetric.
		\end{itemize}
	\item Case $\rt s = 1_X$.
		If $\rb u = 1_Y$ and $\rb v = 1_Z$,
		both sides of the equality are equal to
		$1_{X \msubst Y \msubst Z [y]}$.
		Otherwise, they are equal to
		$\esum{X \msubst Y \msubst Z [y]}$.
	\item Case $\rt s = \rt s' \cdot \rb s''$
		is similar to the case of application above.
	\qedhere
	\end{itemize}
	\end{proof}


\subsection{Strong confluence of \texorpdfstring{$\rred$}{→r}
(Theorem \ref{the:rred-norm-confl}, item \ref{the:rred-norm-confl:confl})}

\begin{lemma} \label{lem:rred-and-subst-1}
	If $\rt s \rred \rs s'$ then 
	$\rt s \rsubst t \rredr \rs s' \rsubst t$.
\end{lemma}
	
	\begin{proof}
	By induction on $\rt s \rred \rs s'$.
	For the base case $(\ensuremath{\lambda}y.\rt u)\rb v \rred \rt u \rsubst v [y]$,
	\begin{align*}
		\left( (\ensuremath{\lambda}y.\rt u)\rb v \right) \rsubst t
		& = \sum_{\rb t = \rb t' \cdot \rb t''}
			\left( \ensuremath{\lambda}y.\rt u \rsubst* {\rb t'} \right)
			\rb v \rsubst* {\rb t''} [y] \\
		& \rredr \sum_{\rb t = \rb t' \cdot \rb t''}
			\rt u \rsubst* {\rb t'}
			\rsubst*{ \rb v \rsubst* {\rb t''} } [y] \\
		& = u \rsubst v [y] \rsubst t
			& \text{by \cref{lem:resource-substitution}.}
	\end{align*}
	Notice that the reflexive closure $\rredr$ of $\rred$
	is necessary to handle the case where the two sums
	are empty.
	The other cases are immediate applications 
	of the induction hypotheses.
	\end{proof}

\begin{lemma} \label{lem:rred-and-subst-2}
	If $\rb t \rred \rbs t'$ then 
	$\rt s \rsubst t \rredr \rt s \rsubst* {\rbs t'}$.
\end{lemma}
	
	\begin{proof}
	By a straightforward induction on the different cases
	defining $\rt s \rsubst t$.
	\end{proof}

\begin{lemma} \label{lem:rred-confl-one-term}
	For all $\rt s \in \rterms**(X)$ and 
	$\rs s', \rs s'' \in \frsums[X]**$ such that 
	\[ \rt s \rred \rs s' \quad \text{and} \quad \rt s \rred \rs s'', \]
	there exists $\rs t \in \frsums[X]**$ such that
	\[ \rs s' \rredr \rs t \quad \text{and} \quad \rs s'' \rredr \rs t. \]
\end{lemma}
	
	\begin{proof}
	Take $\rt s$, $\rs s'$ and $\rs s''$ as above.
	We proceed by induction on both reductions
	$\rt s \rred \rs s'$ and $\rt s \rred \rs s''$.
	The base case is when the first reduction
	comes from the first rule in \cref{def:rred},
	\emph{i.e.} $\rt s = (\ensuremath{\lambda}x.\rt u)\rb v$
	and $\rs s' = \rt u \rsubst v$:
	\begin{itemize}
	\item If the second reduction comes from the same rule,
		then $\rs s'' = \rs s'$ and we set $\rs t \eqdef \rs s'$.
	\item If the second reduction comes from the rule for left application,
		\emph{i.e.} $\rs s'' = (\ensuremath{\lambda}x.\rs u'')\rb v$
		with $\rt u \rred \rs u''$, then
		by \cref{lem:rred-and-subst-1}
		$\rt s' = \rt u \rsubst v \rredr \rs u'' \rsubst v$
		and by \cref{def:rred} $\rt s'' \rredr \rs u'' \rsubst v$.
		(The reflexive closure is needed in the latter case
		because $\rs u''$ may be empty.)
	\item If the second reduction comes from the rule for right application,
		then the proof is as in the previous case,
		using \cref{lem:rred-and-subst-2}
		instead of \cref{lem:rred-and-subst-1}.
	\end{itemize}
	In all other cases, the results immediately follows from
	the induction hypotheses.
	\end{proof}

\setcounter{theoremtmp}{\arabic{theorem}}
\setcounter{theorem}{\arabic{the:rred-norm-confl:confl}}
\begin{theorem}[\cref{the:rred-norm-confl:confl},
recalled from \cpageref{the:rred-norm-confl:confl}]
	For all $\rs s, \rs s', \rs s'' \in \frsums[X]**$ 	such that 
	\[ \rs s \rred \rs s' \quad \text{and} \quad \rs s \rred \rs s'', \]
	there exists $\rs t \in \frsums[X]**$ such that
	\[ \rs s' \rredr \rs t \quad \text{and} \quad \rs s'' \rredr \rs t. \]
\end{theorem}
\setcounter{theorem}{\arabic{theoremtmp}}

	\begin{proof}
	We prove the result under the slightly more general hypothesis
	that $\rs s \rredr \rs s'$ and $\rs s \rredr \rs s''$.
	We proceed by induction on $\sumsize{\rs s}$.
	If $\sumsize{\rs s} = []$, \emph{i.e.} $\rs s$ is empty,
	then both reductions must be equalities
	and the conclusion is trivial.
	Otherwise, suppose $\sumsize{\rs s} \msetgt []$.
	Again, if any of the two reductions
	turns out to be an equality, the result is immediate.
	Otherwise $\rs s \rred \rs s'$ and $\rs s \rred \rs s''$, \emph{i.e.}
	\begin{itemize}
	\item $\rs s = \rt s_1 + \rs s_{\neq 1}$
		and $\rs s' = \rs s'_1 + \rs s'_{\neq 1}$,
		together with $\rt s_1 \rred \rs s'_1$
		and $\rs s_{\neq 1} \rredr \rs s'_{\neq 1}$,
	\item $\rs s = \rt s_2 + \rs s_{\neq 2}$
		and $\rs s'' = \rs s''_2 + \rs s''_{\neq 2}$,
		together with $\rt s_2 \rred \rs s''_2$
		and $\rs s_{\neq 2} \rredr \rs s''_{\neq 2}$.
	\end{itemize}
	There are two possible cases.
	\begin{itemize}
	\item If $\rt s_1 = \rt s_2$ and $\rs s_{\neq 1} = \rs s_{\neq 2}$,
		then by \cref{lem:rred-confl-one-term} applied to
		$\rt s_1 \rred \rs s'_1$ and $\rt s_1 \rred \rs s''_2$,
		we obtain $\rs t_1$ such that
		$\rs s'_1 \rred \rs t_1$ and $\rs s''_2 \rred \rs t_1$.
		By the induction hypothesis on $\rs s_{\neq 1}$,
		we obtain $\rs t_{\neq 1}$ such that
		$\rs s'_{\neq 1} \rred \rs t_{\neq 1}$ 
		and $\rs s''_{\neq 2} \rred \rs t_{\neq 1}$.
		Therefore we can set $\rs t \eqdef \rs t_1 + \rs t_{\neq 1}$.
	\item If $\rt s_1 \neq \rt s_2$, then we can write
		$\rs s = \rt s_1 + \rt s_2 + \rs s_{\neq 12}$, with:
		\begin{align*}
			\rs s_{\neq 1} & = \rt s_2 + \rs s_{\neq 12} &
			\rs s'_{\neq 1} & = \rs s'_2 + \rs s'_{\neq 12} &
			\rt s_2 & \rredr \rs s'_2 &
			\rs s_{\neq 12} & \rredr \rs s'_{\neq 12} \\
			\rs s_{\neq 2} & = \rt s_1 + \rs s_{\neq 12} &
			\rs s''_{\neq 2} & = \rs s''_1 + \rs s''_{\neq 12} &
			\rt s_1 & \rredr \rs s''_1 &
			\rs s_{\neq 12} & \rredr \rs s''_{\neq 12}.
		\end{align*}
		By \cref{lem:rred-confl-one-term} applied to $\rt s_1$
		we obtain $\rs t_1$ such that
		$\rs s'_1 \rredr \rs t_1$ and $\rs s''_1 \rredr \rs t_1$.
		By the same lemma applied to $\rt s_2$
		we obtain $\rs t_2$ such that
		$\rs s'_2 \rredr \rs t_2$ and $\rs s''_2 \rredr \rs t_2$.
		By the induction hypothesis on $\rs s_{\neq 12}$
		we obtain $\rs t_{\neq 12}$ such that
		$\rs s'_{\neq 12} \rredr \rs t_{\neq 12}$
		and $\rs s''_{\neq 12} \rredr \rs t_{\neq 12}$.
		Finally, we can set $\rs t \eqdef 
		\rs t_1 + \rs t_2 + \rs t_{\neq 12}$.
	\qedhere
	\end{itemize}
	\end{proof}

